\section*{Chapter \ref{ch:g12:keynesian-cross} --- The Keynesian Cross and the Multiplier}
% figdata: figdata/grade-12/09-keynesian-cross/out/ (multipliers, brenland_key); Brenland invented.

\begin{solution}{exo:g12:keynesian-cross:1}
$Y = 300 + 0.8Y$: $Y = 300/0.2 = 1500$.
\end{solution}

\begin{solution}{exo:g12:keynesian-cross:2}
$Y = 200 + 0.75Y$: $Y = 200/0.25 = 800$.
\end{solution}

\begin{solution}{exo:g12:keynesian-cross:3}
$1/0.4 = 2.5$ and $1/0.1 = 10$.
\end{solution}

\begin{solution}{exo:g12:keynesian-cross:4}
Four: news shocks (up to 2), Zubairy (1.05), Leeper--Traum--Walker (1.36) and Sims--Wolff (1.07). Coenen et al.\ reach exactly 1.
\end{solution}

\begin{solution}{exo:g12:keynesian-cross:5}
0.3: the spending cuts in 17 OECD countries and the austerity plans in 16, both from budget consolidations, and the lower end of the high-income countries' range.
\end{solution}

\begin{solution}{exo:g12:keynesian-cross:6}
$1/(1 - 0.8 \times 0.75 + 0.1) = 1/0.5 = 2$.
\end{solution}

\begin{solution}{exo:g12:keynesian-cross:7}
$1/(1 - 0.63 + 0.3) = 1/0.67 \approx 1.49$, against $1/(1 - 0.9) = 10$.
\end{solution}

\begin{solution}{exo:g12:keynesian-cross:8}
Output rises by $2 \times 50 = 100$; tax receipts by $0.25 \times 100 = 25$; the deficit by $50 - 25 = 25$.
\end{solution}

\begin{solution}{exo:g12:keynesian-cross:9}
100, 80, 64: total 244 after three rounds, against a limit of $100/0.2 = 500$.
\end{solution}

\begin{solution}{exo:g12:keynesian-cross:10}
1000 and 1500: $500/100 = 5 = 1/(1 - 0.8)$.
\end{solution}

\begin{solution}{exo:g12:keynesian-cross:11}
\omterm{def:g12:keynesian-cross:planned-expenditure}{Autonomous spending} does not change: the multiplier is zero.
\end{solution}

\begin{solution}{exo:g12:keynesian-cross:12}
Almost all the estimates in the chart lie between 0.3 and 1.4, far from the 5, 10 or 20 the cross gives with a high propensity to consume and no other leaks. They also
differ by conditions: below one on average, possibly higher in deep slumps or at the zero lower bound, between 0.1 and $-0.7$ with flexible exchange rates. A multiplier
that is large and the same everywhere fits none of this.
\end{solution}

\begin{solution}{exo:g12:keynesian-cross:13}
The spending rise adds $100/0.2 = 500$; the tax rise lowers consumption by 80 at first, which removes $80/0.2 = 400$. Output rises by 100: a balanced-budget multiplier of 1
in this model.
\end{solution}

\begin{solution}{exo:g12:keynesian-cross:14}
\omterm{def:g12:aggregate-demand-and-supply:ad}{Aggregate demand} shifts right; along the vertical long-run supply curve, output stays at potential and the price level rises. The multiplier is zero in the long run.
\end{solution}

\begin{solution}{exo:g12:keynesian-cross:15}
Imports: part of each round of extra spending goes to foreign producers, so less returns as domestic income. Floating rate: if the central bank raises its rate, the
currency rises and \omterm{def:g10:national-accounts:net-exports}{net exports} fall, an extra offset; under a fixed rate the central bank keeps its rate in line with the anchor country. The debate's estimates: about
1.4 in the long run with fixed rates, between 0.1 and $-0.7$ with flexible ones.
\end{solution}

\begin{solution}{pb:g12:keynesian-cross:1}
\textbf{1.} $393\,264 - 157\,306 = 235\,958$ million; $225\,014/235\,958 \approx 0.954$.

\textbf{2.} $157\,306/393\,264 \approx 0.400$.

\textbf{3.} $1/(1 - 0.95) = 20$.

\textbf{4.} $0.01 \times 393\,264 \approx 3933$ million dalens; $20 \times 3933 = 78\,660$ million, 20\,\% of GDP.

\textbf{5.} $1/(1 - 0.95 \times 0.6) = 1/0.43 \approx 2.33$.

\textbf{6.} $1/(1 - 0.57 + 0.34) = 1/0.77 \approx 1.30$.

\textbf{7.} $3933/0.77 \approx 5108$ million dalens.

\textbf{8.} $0.40 \times 5108 \approx 2043$ million; the deficit rises by $3933 - 2043 = 1890$ million.

\textbf{9.} $0.34 \times 5108 \approx 1737$ million.

\textbf{10.} It assumes that firms can raise output by a fifth with no rise in prices, at an unchanged interest rate, and that taxes and imports take nothing from each
round, none of which holds in Brenland.

\textbf{11.} \omterm{def:g12:keynesian-cross:planned-expenditure}{Autonomous spending} rises by $3933/2 \approx 1967$ million: a gain of $1966.5/0.77 \approx 2554$ million.

\textbf{12.} $2554/3933 \approx 0.65$.

\textbf{13.} Zero: the price level would rise instead of output.

\textbf{14.} 1.30 lies near the top of the chart's estimates; 0.65 lies inside the range of most of them.

\textbf{15.} That households spend a fixed share of current income and do not look ahead.

\textbf{16.} $1/(1 - 0.57 + 0.17) = 1/0.6 \approx 1.67$.

\textbf{17.} Its central bank keeps its rate tied to the anchor currency and does not raise it in response, so the monetary offset is absent.

\textbf{18.} A one-off rise adds little to their income over the years; households that look ahead spread it out or save it, so $c$ for a temporary rise is lower.

\textbf{19.} The one with all three offsets, 0.65, with the warning that the size of the central bank's response is assumed, not estimated, and that the toy model has
fixed prices.

\textbf{20.} The toy-model multiplier, with consumption as the only leak, is 20; counting taxes and imports leaves 1.30, and the monetary offset leaves 0.65.
\end{solution}
